% =====================================================================
% SkillRACE — main paper. Style mirrors the lab's gordian.tex (acmart,
% nonacm). Uses `listings` (not minted) so it compiles without
% -shell-escape. All not-yet-measured quantities are \newcommand macros
% grouped in the PLACEHOLDERS block below: fill them once, they update
% everywhere, and every occurrence is visibly a placeholder until then.
% =====================================================================
\documentclass[acmsmall,screen,nonacm]{acmart}

\usepackage{booktabs}
\usepackage{listings}
\usepackage{xcolor}
\usepackage{graphicx}
\usepackage{algorithm2e}
\usepackage{tikz}
\usetikzlibrary{arrows.meta,positioning,calc,fit,backgrounds}
\usepackage{makecell}
\usepackage{caption}
\usepackage{multirow}
\usepackage{array}
\usepackage{cleveref}
\usepackage{subcaption}
\usepackage[most]{tcolorbox}
\tcbuselibrary{breakable}
\usepackage{enumitem}
\setlist{leftmargin=6mm}

% ---- revision toggle (as in gordian) ----
\newif\ifrevision
\revisionfalse
\ifrevision \newcommand{\rev}[1]{\textcolor{red}{#1}}
\else \newcommand{\rev}[1]{#1}\fi

% ---- author-note macros ----
\newcommand{\todo}[1]{{\color{purple}\textbf{[TODO: #1]}}}
\newcommand{\ph}[1]{{\color{blue}#1}} % renders a placeholder value in blue

% ---- project + benchmark names ----
\newcommand{\tool}{\textsc{SkillRACE}\xspace}

% ---- highlight boxes (as in gordian) ----
\newtcbox{\hly}{on line,arc=0pt,colframe=yellow!50,colback=yellow!50,
  boxrule=0pt,boxsep=0pt,left=1pt,right=1pt,top=2pt,bottom=2pt}
\newtcbox{\hlg}{on line,arc=0pt,colframe=green!45,colback=green!45,
  boxrule=0pt,boxsep=0pt,left=1pt,right=1pt,top=2pt,bottom=2pt}
\newtcbox{\hlgr}{on line,arc=0pt,colframe=gray!25,colback=gray!25,
  boxrule=0pt,boxsep=0pt,left=1pt,right=1pt,top=2pt,bottom=2pt}

% ---- listing style ----
\lstdefinestyle{skill}{
  basicstyle=\ttfamily\scriptsize,
  keywordstyle=\color{blue!70!black},
  commentstyle=\color{gray},
  breaklines=true,
  columns=fullflexible,
  frame=single,
  framesep=3pt,
  xleftmargin=4pt,xrightmargin=4pt,
  tabsize=2,
}

% =====================================================================
% PLACEHOLDERS — fill these once experiments land.
% =====================================================================
\newcommand{\NumSkills}{\ph{28}}            % skills in the suite (24 built + 4 build-deferred)
\newcommand{\NumSkillsCore}{\ph{TBA}}       % skills w/ full property matrix
% D2 (skill-generation) suite — concrete
\newcommand{\NumScenarios}{10}
\newcommand{\NumTestsPer}{10}
\newcommand{\NumTestsTot}{100}
\newcommand{\NumChecksTot}{192}
\newcommand{\Budget}{\ph{120}}              % agent runs per skill per method
\newcommand{\SeedCount}{\ph{20}}            % seed inputs per campaign
\newcommand{\AgentModel}{\ph{a fixed frontier model}}
\newcommand{\JudgeModel}{\ph{a fixed cheap model}}
% headline bug-finding deltas (percentage points / factors) — TBA
\newcommand{\BugsSkillrace}{\ph{TBA}}
\newcommand{\BugsGreybox}{\ph{TBA}}
\newcommand{\BugsRandom}{\ph{TBA}}
\newcommand{\DeltaOverGreybox}{\ph{TBA}}
\newcommand{\DeltaOverRandom}{\ph{TBA}}
% injection detection rate
\newcommand{\DetectRate}{\ph{TBA}}
% claim-2 hidden-test pass rates
\newcommand{\PassZeroShot}{\ph{TBA}}
\newcommand{\PassRandom}{\ph{TBA}}
\newcommand{\PassGreybox}{\ph{TBA}}
\newcommand{\PassSkillrace}{\ph{TBA}}
% calibration agreement
\newcommand{\SegF}{\ph{TBA}}
\newcommand{\MergeAgree}{\ph{TBA}}
\newcommand{\CompileAgree}{\ph{TBA}}

% =====================================================================
\title[\tool: Reasoning-Augmented Concolic Execution for Coding-Agent Skills]{%
  \tool: Reasoning-Augmented Concolic Execution\\for Coding-Agent Skills}

\author{Anonymous}
\affiliation{\institution{Anonymous Institution}\country{Anonymous}}
\email{anon@example.com}

\begin{document}

\begin{abstract}
A coding-agent \emph{skill} is reusable procedural guidance---a natural-language
\texttt{SKILL.md} plus optional scripts---that an agent consults to carry out a
task such as fixing a failing test or scaffolding a service. Skills are now a
first-class, sharable artifact, yet they are trusted after a handful of successful
runs and tested almost not at all: unlike code, a skill has no compiler and no test
suite, and the guidance that works on the situations its author happened to try may
mislead an agent on the situations its author did not. We present \tool, which tests
a skill the way concolic execution tests code. \tool runs the skill, lifts each run
into a sequence of \emph{episodes}---coherent units of work, each with a purpose and
an outcome read from tool output rather than from the agent's own narration---and
merges episodes into a growing \emph{behavior tree}. It reads the agent's reasoning
at each branch as the \emph{condition} that sent the run one way rather than another,
then mutates those conditions to synthesize new \((\text{prompt},\text{environment})\)
inputs that drive the skill down branches it has not yet taken. Candidate inputs are
validated against the target condition \emph{before} any agent run, so the expensive
resource---the agent under test---is spent only on inputs already shown to exercise
the intended situation. Correctness is judged by \emph{properties}: universal
process invariants checked mechanically, and per-skill specifications compiled into
executable checks \emph{before} the run they judge, so that no model ever grades the
run it also produced. We evaluate \tool against a blind-mutation floor and a
greybox rung that ports the feedback mechanism of a recent agent fuzzer to this
setting, on a suite of \NumSkills\ skills mined from public repositories. Holding the
runner, environments, budget, and oracle identical across methods, \tool exposes
\DeltaOverGreybox\ more distinct defects per agent-run budget than the greybox rung and
\DeltaOverRandom\ more than the floor, and reaches its first violation in fewer runs.
Feeding \tool's findings back to revise an LLM-generated skill raises its pass rate on
independently-authored held-out tests from \PassZeroShot\ to \PassSkillrace.
\end{abstract}

\maketitle

% =====================================================================
\section{Introduction}
\label{sec:intro}

Coding agents increasingly rely on \emph{skills}: reusable folders of natural-language
guidance, and sometimes helper scripts, that an agent loads on demand to perform a
recurring task~\cite{anthropic2025skills,anthropic2025skillseng}. A skill for fixing
a failing test tells the agent to locate the failure, understand the expected
behavior, edit the implementation, and re-run the suite; a skill for scaffolding a
service prescribes a project layout and a build step. Skills are authored once,
published to shared directories, and reused across many tasks and many
agents~\cite{anthropic2025skillseng}. In this they resemble a library---code that
many downstream users depend on---except that a skill is prose, its ``interface'' is
whatever situations its author imagined, and it ships with neither a type checker nor
a test suite.

This creates a testing gap that is easy to state and hard to close. A skill is
trusted once it has worked on a handful of tasks its author tried by hand. But the
tasks an author tries are the tasks an author thinks of, and a skill's guidance is
only as good as its coverage of the situations an agent will actually meet. Guidance
that is silent about a broken build, ambiguous about which of two files to edit, or
that assumes a test harness runs cleanly will send an agent confidently down the
wrong path exactly when the environment departs from the author's mental model---and
because the agent narrates its own progress, the failure is often invisible in the
transcript the author reads. What is missing is a way to ask, systematically, \emph{on
which situations does this skill's guidance break down?}

The analogous question for code has a mature answer. Symbolic and concolic execution
run a program on abstract inputs, read the branch conditions off the source, and negate
them to synthesize inputs that drive execution down not-yet-taken
paths~\cite{king1976symbolic,dart,sen2005cute,klee}. Coverage of the path space is
both the goal and the search signal: the path tree built so far reveals the next
frontier to push toward. This is precisely the discipline a skill lacks---and it does
not transfer directly, because a skill is not code. There is no source to read a
branch condition from; the ``program'' is an agent whose control flow is mediated by
natural-language reasoning over the situations it observes.

Our key observation is that the agent's reasoning trace is the nearest thing to that
missing source. When an agent writes ``the import failed, so the module is missing---I
will create it'' before choosing its next action, it is stating the condition under
which it took that action rather than another. \tool treats each such stated reason as
a candidate branch condition. It segments a run into \emph{episodes}---contiguous
stretches of work pursuing one sub-goal, the basic blocks of agent behavior---and
merges episodes across runs into a \emph{behavior tree} whose branches are the points
where runs that reached the same situation diverged. At each branch it reads two
signals: the \emph{outcome} of the episode that just finished (taken from tool output,
not from the agent's claim about it) and the \emph{opening reasoning} of the episode
that follows. Mutating that condition---negating a binary guard, or targeting a
sibling value of a multi-valued one---yields a new \((\text{prompt},\text{environment})\)
input designed to drive the skill down an unexplored branch. Each new run folds back
into the tree, exposing the next frontier, exactly as in concolic execution.

Two design commitments make this practical rather than merely appealing. First, the
agent under test is the only expensive component; everything else is code or a small
model. \tool therefore never spends an agent run on an unvalidated input: a synthesized
candidate is built into its container and checked against its target condition---does
the test it sets up actually fail with the intended error?---with no agent involved,
and only validated candidates are run. Second, correctness is decided by
\emph{properties}, not by a model's opinion of the run. Universal process invariants
(a skill must not force-push, must not delete the very test it was asked to satisfy,
must terminate within budget) are checked mechanically. Skill-specific expectations
are written once as natural-language specifications and \emph{compiled per input into
an executable check before the run they judge}, so the component that decides whether
a run is buggy never sees the run it grades. A run is buggy when it violates a property
we fixed in advance---not when a model thought it went badly.

We use \tool to answer two questions. \emph{Does reasoning-guided exploration find
more skill defects than feedback that ignores reasoning?} We compare against a
blind-mutation floor and a greybox rung that ports the tool-invocation-sequence
feedback of a recent agent fuzzer~\cite{verigrey} to correctness testing, holding the
runner, the environments, the budget, and the property checker identical across all
three so that any difference measures test \emph{generation}, not detection.
\emph{Can a skill be improved with what \tool finds?} Most skills in circulation are
themselves LLM-generated, which is circular---a model writing down what it does or does
not reliably know how to do. We test whether \tool breaks the circle: we revise an
LLM-generated skill using \tool's reports and measure the revised skill on held-out
tests authored independently of the revision loop.

\smallskip\noindent This paper makes the following contributions:
\begin{itemize}[itemsep=2pt,topsep=2pt]
  \item \textbf{A concolic-execution account of agent-behavior testing} that lifts the
  branch-condition/negation loop from source code to the agent's reasoning trace:
  episodes as basic blocks, stated reasons as branch conditions, condition mutation as
  test synthesis, and a behavior tree as both the byproduct and the driver of the
  search (\cref{sec:overview,sec:approach}).
  \item \textbf{A cost-aware realization} of that loop in which a cheap validator gates
  every expensive agent run, and correctness is judged by fixed invariants plus
  specifications compiled into checks \emph{before} the run they grade
  (\cref{sec:approach,sec:impl}).
  \item \textbf{An evaluation} on \NumSkills\ skills mined from public repositories,
  against a blind-mutation floor and a faithful greybox port, measuring defect yield
  under a fixed budget; plus a study of whether \tool's findings improve an
  LLM-generated skill on held-out tests (\cref{sec:eval}).
\end{itemize}

% =====================================================================
\section{Overview}
\label{sec:overview}

\begin{figure*}[t]
  \centering
  \resizebox{\linewidth}{!}{%
  \begin{tikzpicture}[
    font=\footnotesize\sffamily,
    >={Stealth[length=1.7mm,width=1.6mm]},
    node distance=9mm and 12mm,
    box/.style={rounded corners=2.5pt, draw=black!45, line width=0.5pt,
                minimum height=11mm, minimum width=23mm, align=center,
                inner sep=3pt, fill=white},
    model/.style={box, fill=black!5},                      % cheap-model step
    code/.style={box, fill=white},                          % code-only step
    agent/.style={box, fill=orange!14, draw=orange!70!black, line width=1pt}, % expensive
    art/.style={box, fill=blue!7, draw=blue!45!black, densely dashed},        % artifact
    io/.style={align=center, inner sep=1pt, text=black!65, font=\scriptsize\sffamily},
    flow/.style={->, draw=black!55, line width=0.7pt},
    gate/.style={->, draw=orange!75!black, line width=1pt},
    lbl/.style={font=\scriptsize\sffamily, text=black!60, inner sep=1.5pt},
  ]
    % --- observe path (top row, left to right) ---
    \node[agent] (run)  {\textbf{Runner}\\[-1pt]{\scriptsize agent under test}};
    \node[model] (seg)  [right=of run]  {Segmenter\\[-1pt]{\scriptsize $\to$ episodes}};
    \node[art]   (tree) [right=of seg]  {Behavior\\[-1pt]tree};
    % --- generate path (bottom row, right to left) ---
    \node[model] (guard)[below=of tree] {Guard extractor\\[-1pt]{\scriptsize + property-guided}\\[-1pt]{\scriptsize selection}};
    \node[model] (synth)[left=of guard] {Synthesizer};
    \node[code]  (valid)[left=of synth] {Validator\\[-1pt]{\scriptsize no agent}};
    % --- inputs / oracle ---
    \node[io]    (seed) [left=8mm of run]  {seed\\inputs};
    \node[code]  (check)[above=of run]     {Property checker\\[-1pt]{\scriptsize fixed invariants +}\\[-1pt]{\scriptsize pre-compiled checks}};
    \node[io]    (bugs) [left=8mm of check]{bug\\reports};
    % --- flows ---
    \draw[flow] (seed) -- (run);
    \draw[flow] (run)  -- (seg)  node[lbl,midway,above]{trace};
    \draw[flow] (seg)  -- (tree);
    \draw[flow] (tree) -- (guard) node[lbl,midway,right]{frontier};
    \draw[flow] (guard)-- (synth);
    \draw[flow] (synth)-- (valid) node[lbl,midway,above]{candidate};
    \draw[gate] (valid) -- (run)  node[lbl,midway,left,text=orange!70!black,align=right]{validated\\input only};
    \draw[flow] (run)  -- (check) node[lbl,midway,right,align=left]{each run's\\final state};
    \draw[flow] (check)-- (bugs);
    % --- container behind the loop, labelled ---
    \begin{scope}[on background layer]
      \node[rounded corners=4pt, fill=black!3, draw=black!12,
            fit=(run)(seg)(tree)(guard)(synth)(valid),
            inner sep=6mm, label={[lbl,anchor=north east]north east:exploration loop}] {};
    \end{scope}
    % --- legend, stacked in the whitespace to the right of the loop ---
    \begin{scope}[every node/.append style={anchor=west},
                  swatch/.style={minimum width=5mm, minimum height=3.4mm, inner sep=0pt}]
      \coordinate (lg) at ($(tree.east)+(14mm,4mm)$);
      \node[agent, swatch] (le1) at (lg) {};
      \node[io] at (le1.east) {\,expensive (agent)};
      \node[model, swatch] (le2) at ($(le1.south west)+(0,-2.4mm)$) {};
      \node[io] at (le2.east) {\,cheap model};
      \node[code, swatch] (le3) at ($(le2.south west)+(0,-2.4mm)$) {};
      \node[io] at (le3.east) {\,code only};
      \node[art, swatch] (le4) at ($(le3.south west)+(0,-2.4mm)$) {};
      \node[io] at (le4.east) {\,artifact};
    \end{scope}
  \end{tikzpicture}}
  \caption{\tool's loop. A run is segmented into episodes and folded into the behavior
  tree; the tree's branch frontier drives property-guided synthesis of a new input,
  which the \emph{validator} confirms exercises the targeted branch \emph{before} any
  agent run. The agent under test (orange) is the only expensive stage and is entered
  only through the validated-input gate; every other stage is a cheap model call or
  plain code. Each run's final container is checked against fixed invariants and the
  input's pre-compiled checks, yielding bug reports.}
  \label{fig:pipeline}
\end{figure*}

We illustrate \tool with a defect it found in a \texttt{fix-failing-test} skill---a
common skill whose guidance is, in essence, ``run the suite, read the failing tests,
fix the implementation, re-run to confirm.'' The defect is not a flaw in any single
run; it is a situation the guidance does not cover, surfaced during a \tool campaign
and then generalized by the search into a dimension worth probing.

\paragraph{A run that looks fine and is not.} \tool seeds the campaign with inputs
that place a bug in a project and ask the agent to make the tests pass. One seed
placed the bug in a module the setup happened to name \texttt{collections.py}. This
name shadows a Python standard-library module, so the test harness crashes on
\emph{import}, before any test runs---\texttt{pytest} itself fails with
\texttt{ImportError}. The skill's guidance says nothing about a harness that will not
start. The agent edited the implementation, could not get the suite to run, and
eventually ``verified'' its fix by executing the functions directly, outside the test
runner---then reported success. Its narration reads like a clean repair. The tool
output tells a different story: the suite never passed, and no successful test run ever
occurred.

\paragraph{Why \tool catches it.} \tool does not ask a model whether the run went
well. It reads the outcome of each episode from tool output. The episode that ran the
suite has outcome ``\texttt{ImportError} on import; no tests collected,'' regardless of
what the agent said next. Two properties compiled for this input from the skill's
specification---\emph{the previously-failing suite passes in the final state} and
\emph{the agent actually ran the suite to a clean result before finishing}---are
checked by scripts that grep the final workspace and the trace. Both fire. The report
names the violated properties, the input that produced them, and a replayable
container.

\paragraph{From one finding to a systematic probe.} The value of \tool is not the single
run above but the \emph{dimension} it exposes. Folding the campaign's runs into the
behavior tree produced a branch on \emph{which module the workspace directs the agent to
fix}; \tool extracted this as a guard, grounded it in an executable check over the initial
state (\texttt{test -f <module>}), and, asked which feasible mutation most threatened a
property, synthesized further inputs that vary the targeted module and its failure mode.
Module identity turned out not to change the agent's behavior on its own (the loop records
this as a non-divergence and moves on), whereas the \emph{shape of the harness failure}---a
suite that will not even start---does. That distinction is exactly what a reasoning-blind
fuzzer cannot draw and what \tool's guards make explicit: the search is steered by the
conditions a skill's own runs reveal and the properties those conditions threaten, not by
random perturbation. \todo{once the full campaign lands, replace with a synthesized input
that itself triggered a violation, if one is cleaner than this seed-found case.}

% =====================================================================
\section{Approach}
\label{sec:approach}

\tool is a loop over six components (\cref{fig:pipeline}). One---running the agent under
test---is expensive: minutes and real cost per run. The other five are code or a small
model. Every design choice follows from this asymmetry: \emph{spend agent runs only on
inputs already validated, and do everything else cheaply.} We first fix the objects the
components exchange, then describe each.

\subsection{Runs, episodes, and the behavior tree}
\label{sec:objects}

\paragraph{Run.} The runner executes the skill's agent in a container on an input
\((x, E_0)\): a prompt \(x\) and an initialization that builds the starting environment
\(E_0\). It records the full trace---for each step, the agent's reasoning, the tool call,
and the raw tool output---and leaves the final container live for property checking.
Coding skills are short, so traces are tens of steps, not thousands.

\paragraph{Episode.} An \emph{episode} is a contiguous run of steps pursuing one
sub-goal: ``reproduce the failure,'' ``edit the implementation,'' ``run the suite.'' It
is the basic block of the behavior tree: internally it may be one tool call or eight,
but it is treated as a unit with an \emph{intent} (its sub-goal) and an \emph{outcome}
(its result). Boundaries are read from the reasoning, because the agent announces its
own goal switches. Crucially, the outcome is read from the episode's \emph{tool output}
(an exit code, an error string, an HTTP status), never from the agent's narration of how
it went. Agents declare false victory; sourcing the outcome from observation is what let
\tool see through the run in \cref{sec:overview}, and---when the agent's next reasoning
proceeds as though the outcome were different from what the tool reported---the
disagreement is itself a defect signal.

\paragraph{Behavior tree.} Runs are merged into a tree whose nodes are situations the
skill reaches and whose edges are what it did next. A new run is folded from the root:
at each node, its next episode is compared, \emph{on purpose}, against the node's
children. A match descends into that child and broadens the child's purpose to cover
the new episode; a mismatch creates a new child. Matching is on purpose, not on
outcome: the outcome rides on the \emph{edge}. So two runs that reach the same situation
but fare differently share the node and carry different outcomes on their out-edges---and
the divergence, if it matters, surfaces one step later as a \emph{branch}, a node with
two or more children. A branch is where a guard lives.

\subsection{Guards: reading a branch condition from reasoning}
\label{sec:guards}

A guard is the condition that explains a branch. \tool reads it from two signals with
distinct sources, which it is careful not to conflate:
\begin{itemize}[itemsep=1pt,topsep=2pt]
  \item the \textbf{outcome} of the episode that just finished, from its tool output---%
  \emph{what state the agent was in}; and
  \item the \textbf{opening reasoning} of the next episode, from the agent's reasoning
  text---\emph{why it went the way it did, given that state}.
\end{itemize}
A model call takes these two signals for each diverging side, plus a diff of the sides'
initial environments, and distills a condition together with, where possible, an
\emph{executable grounding}: a concrete check such as \texttt{test -f <path>} or
``\texttt{pytest} output contains \texttt{ImportError}.'' It also records
\emph{decidability}: a guard is targetable now if its condition is decidable from the
initial setup \(E_0\) (``the module is missing,'' ``the failing test is an import
error''); guards decidable only by running the agent are deferred and counted, not
targeted. This is the conceptual core: a program hands the tester its branch conditions
in source; an agent's reasoning is the nearest surrogate, and \tool treats each stated
reason as a \emph{candidate} condition to be confirmed or refuted by mutation.

\subsection{Property-guided synthesis, validated before any agent run}
\label{sec:synth}

From the tree's branches \tool builds a \emph{frontier}: the untried mutations of
\(E_0\)-decidable guards. To explore, it must choose one---and the choice is where the
search earns its keep. \tool shows the selection step the skill's \emph{properties}
alongside the frontier and asks for the feasible mutation most likely to drive the skill
toward violating one: if the guard took this other value, could the guidance plausibly
mishandle it in a way a property would catch? This turns exploration into directed
falsification---the frontier supplies what is reachable, the properties supply what is
worth reaching---while leaving the verdict to the property checker, not the selector.

A cheap generator then drafts a candidate input for the chosen mutation: a prompt and an
environment intended to reach the branch \emph{and} realize the mutated condition,
together with a check script for that condition. The generator may be unreliable; the
\emph{validator} is not. It builds the candidate's container and runs the check with no
agent involved---does the test it set up actually fail with the targeted error?---and a
botched candidate costs a rebuild, never an agent run. Only a validated input is handed
to the runner. We summarize the loop in \cref{alg:loop}.

\begin{algorithm}[t]
\SetAlgoLined\DontPrintSemicolon
\footnotesize
\KwIn{skill $S$, properties $P$, seed inputs, budget $B$}
\KwOut{behavior tree $T$, bug reports $R$}
$T \leftarrow \emptyset$;\quad $R \leftarrow \emptyset$\;
\ForEach{seed input $(x,E_0)$}{
  run agent; check $P$; fold run into $T$ \tcp*{seed phase}
}
\While{agent runs spent $< B$}{
  $F \leftarrow$ frontier of $T$ \tcp*{untried $E_0$-decidable mutations}
  $(g, m) \leftarrow$ \textsc{Select}$(F, P)$ \tcp*{property-guided}
  $(x^*,E_0^*, c) \leftarrow$ \textsc{Synthesize}$(g, m)$ \tcp*{+ check script $c$}
  \If{\textbf{not} \textsc{Validate}$(E_0^*, c)$}{mark $m$ tried; \textbf{continue}\tcp*{no agent spent}}
  run agent on $(x^*,E_0^*)$\;
  check $P$; add violations to $R$\;
  fold run into $T$; classify vs.\ targeted branch\;
}
\Return $T$, $R$\;
\caption{\tool exploration loop.}
\label{alg:loop}
\end{algorithm}

Each new run is folded back and classified against the branch it targeted:
\emph{predicted divergence} (it took the new branch---coverage gained, the frontier
extends), \emph{no divergence} (it reached the branch but behaved as before---the stated
reason was not causal for this mutation, which is marked spent), or \emph{path miss} (it
never reached the branch---an earlier guard failed). No-divergence is not a failure of
the tool; it is the loop confirming that a hypothesized condition did not matter, which
is information the search uses.

\subsection{Properties: the oracle}
\label{sec:props}

A run is buggy when it violates a property fixed in advance. Properties come in two
kinds, reported separately by trustworthiness.

\emph{Fixed invariants} are universal, task-independent, and checked by plain code over
the trace and final state: no force-push; no destructive operation outside the
workspace; no pathological repetition; termination within budget; and---the
reward-hacking core---never satisfy a target test by editing or deleting the test rather
than the code. These involve no model at all, at authoring or evaluation time.

\emph{Specification-by-example} properties capture what a specific skill must achieve---%
``the previously-failing suite passes,'' ``the built tool performs the requested
function.'' They are written once per skill in natural language and \emph{compiled per
input into an executable check}. The compiling model sees the prompt and the built
initial environment---its tools, its starting files---but \textbf{never the run it will
judge}; the produced check is a bash script that discovers concrete targets
mechanically (find the test files, grep the trace) and decides the verdict by exit code.
Because the check attaches to the input, every run of that input---under any method---is
judged by byte-identical scripts, which is also a fairness control. The model runs only
at compile time; evaluation is deterministic and auditable, and the generated scripts
ship with the reports.

This design keeps a bright line: no model ever grades the run it also produced. Bug
yield is of course bounded by property strength---a defect no property forbids is
invisible, as in all assertion-based testing---so we do not claim completeness. We claim
the converse, and measure it: a violation is a real defect, and we quantify how many
planted defects the property suite detects (\cref{sec:eval}).

% =====================================================================
\section{Implementation}
\label{sec:impl}

\tool is implemented in Python. The agent under test runs on the Pi agent
harness~\cite{pi2025} inside Docker; \tool interacts with it only through its recorded
session trace and the live final container, so it is agnostic to the agent's internals.
A single model---\JudgeModel---performs every model-driven \tool step (segmentation,
merge judgments, guard extraction, selection, synthesis, check compilation), decoupled
from the agent under test (\AgentModel). We do not mix models across steps; instead we
ablate the choice of \tool's model and report the pipeline's sensitivity to it. All
model-driven steps run at temperature~0 and cache their judgments, so a campaign is a
deterministic procedure given its seeds and selection order.

\paragraph{Trace and episodes.} A deterministic renderer projects Pi's session log into
a globally-numbered list of tool calls with reasoning inline; a boundary may fall only at
a call that carries reasoning, so every episode edge is a real piece of reasoning. The
segmenter emits episode spans validated to partition the calls in order, with one repair
attempt on a malformed split before the trace is flagged unsegmentable and counted.

\paragraph{Tree and guards.} Merge decisions and purpose-broadening are content-hash
cached, so the same episode pair gets the same verdict within a campaign. Guard state
lives beside the tree; deferred (non-\(E_0\)) guards are recorded so the fraction of the
branch space we cannot yet target is reported rather than hidden.

\paragraph{Property checking.} Fixed invariants run on the host in Python. Compiled
checks run as bash inside the live final container, which also exposes the trace and the
diff, so a single mechanism covers both state and trace properties. A syntactically
broken check is re-authored once with the error fed back and, failing that, recorded as
inconclusive---never as a violation.

% =====================================================================
\section{Evaluation}
\label{sec:eval}

We ask:
\begin{description}[leftmargin=3.2em,itemsep=1pt]
  \item[RQ1] Does \tool find more skill defects, under a fixed agent-run budget, than a
  blind-mutation floor and a greybox rung that uses tool-sequence feedback but no
  reasoning?
  \item[RQ2] How do \tool's model-driven components (segmentation, merging, check
  compilation) agree with human judgment, and how strong is the property suite?
  \item[RQ3] Does revising an LLM-generated skill with \tool's findings improve its pass
  rate on held-out tests, relative to revising with the baselines' findings?
  \item[RQ4] How does \tool's yield depend on its model choice and on
  property-guided selection (ablation)?
\end{description}

\subsection{Setup}
\label{sec:setup}

\paragraph{Skill suite.} We evaluate on \NumSkills\ skills mined from public
repositories and skill marketplaces under a \emph{pre-registered} selection protocol
(committed before any campaign) that fixes the sources and the inclusion/exclusion
criteria, logs every admit/reject decision with its reason, and commits to per-skill
reporting---so skill choice cannot be influenced by results. A skill is admitted only if
it (i) produces a \emph{code} artifact whose behavior is mechanically checkable (a suite
that passes, a CLI that computes a value, a query that returns the right answer, a build
that succeeds), (ii) builds in a credential-free container, and (iii) is procedural
rather than a one-liner. Crucially, \emph{environment-contingency}---whether the skill's
behavior depends on the starting state, which is where reasoning-guided exploration can
separate from reasoning-blind feedback---is \emph{recorded but not used for inclusion};
excluding low-contingency skills would bias the comparison in our favor. Document-%
generation skills (creating a \texttt{.docx}, a slide deck) are excluded: their success
is presentational, so no trustworthy mechanical oracle exists. Each admitted skill ships
a base image, a natural-language property specification, and a per-skill applicability
matrix. A mining note worth reporting: the marketplace yields far fewer \emph{distinct}
checkable skills than raw counts suggest---most entries are forks of a few families
(test-driven development, debugging, refactoring), tool-coupled, or presentational---so
we assemble the suite as multiple independent community implementations per family.
\Cref{tab:skills} summarizes it.

\begin{table}[t]
\centering\small
\caption{Skill suite (excerpt of \NumSkills; families assembled as multiple community
implementations each).}
\label{tab:skills}
\begin{tabular}{@{}l l c c@{}}
\toprule
\textbf{Skill} & \textbf{Artifact checked} & \makecell{\textbf{\#props}} & \makecell{\textbf{env-}\\\textbf{contingent}}\\
\midrule
\texttt{fix-failing-test}   & suite passes; tests unedited      & 4 & \checkmark \\
\texttt{build-python-cli}   & CLI runs; produces output         & 5 & \checkmark \\
\texttt{mcp-server-patterns}& project builds; server responds  & 6 & \checkmark \\
\texttt{test-driven-development}& suite passes; test-first in trace & 4 & \checkmark \\
\texttt{sql-queries}        & query returns correct answer      & 3 & \checkmark \\
\texttt{regex-expert}       & validator matches spec; anchored  & 3 & $\circ$ \\
\texttt{frontend-design}    & renderable artifact; responsive   & 8 & $\circ$ \\
\multicolumn{4}{c}{\ph{$\cdots$ \NumSkills\ skills total (cli, refactor, debugging, parser, config, \dots) $\cdots$}}\\
\bottomrule
\end{tabular}
\end{table}

\paragraph{Baselines: a three-rung ladder.} All rungs share \tool's runner, environments,
budget (\Budget\ agent runs per skill), the shared realize/build pipeline, and---%
critically---the \emph{same property checker with the same per-input compiled checks}, so
any difference measures test generation, not detection. All rungs also start from the same
\SeedCount\ seed inputs, produced by proposing diverse (task, environment) ideas from the
skill's specification and building each into a container; the methods differ only in what
they generate \emph{after} the seed phase.
\emph{Floor} mutates seed prompts and environments at random with no behavioral feedback.
\emph{Greybox} ports the feedback mechanism of VeriGrey~\cite{verigrey}, a recent agent
fuzzer: it keeps a run's candidate as a seed when the run's tool-invocation sequence is
novel (a new tool, a new transition, or a new sequence) and assigns exploration energy
by that novelty---but over correctness testing rather than injection, with VeriGrey's
injection-specific mutation and oracle replaced by the shared components. Because raw
tool arguments are near-unique in a coding agent (which would make every run ``novel''
and collapse the feedback to blind mutation), we schematize tool events at three
granularities and, to avoid handicapping the baseline, report the sensitivity and use the
\emph{best-performing} granularity for the headline comparison. \emph{\tool} is the full
loop of \cref{sec:approach}. The Floor$\to$Greybox gap isolates the value of any
behavioral feedback; the Greybox$\to$\tool gap---the headline---isolates the value of the
episode abstraction and reasoning-guided mutation over reasoning-blind novelty.

\paragraph{Metrics.} We report, per skill and pooled: distinct properties violated within
the budget; runs and wall-clock to first violation; and unique behavior-tree branches
covered. A flagged violation is re-run \(k{=}3\) times to separate a reproducible defect
(3/3) from brittleness. We report yield split by property provenance (fixed vs.\
compiled), the compiled ones carrying a wider confidence band.

\subsection{RQ1: Defect yield}
\label{sec:rq1}

\begin{table}[t]
\centering\small
\caption{Defect yield under a fixed budget of \Budget\ agent runs per skill. Values are
placeholders pending the full campaign. \todo{fill.}}
\label{tab:rq1}
\begin{tabular}{@{}l c c c@{}}
\toprule
& \textbf{Floor} & \textbf{Greybox} & \textbf{\tool}\\
\midrule
Distinct properties violated (pooled) & \BugsRandom & \BugsGreybox & \BugsSkillrace\\
Reproducible ($3/3$) violations       & \ph{TBA} & \ph{TBA} & \ph{TBA}\\
Runs to first violation (median)      & \ph{TBA} & \ph{TBA} & \ph{TBA}\\
Unique branches covered               & \ph{TBA} & \ph{TBA} & \ph{TBA}\\
\bottomrule
\end{tabular}
\end{table}

\begin{figure}[t]
\centering
\fbox{\parbox{0.9\columnwidth}{\centering\vspace{15mm}
\textbf{[Figure~2 placeholder]}\\[1mm]
distinct properties violated (y) vs.\ agent runs spent (x),\\
one curve per method, pooled over skills; \tool above Greybox above Floor.
\vspace{15mm}}}
\caption{Defect discovery vs.\ budget. A per-run advantage for \tool would show as a
curve that rises faster and higher. \todo{plot from campaign logs.}}
\label{fig:yield}
\end{figure}

\Cref{tab:rq1} and \cref{fig:yield} will report yield. We expect \tool to exceed Greybox
by \DeltaOverGreybox\ and Floor by \DeltaOverRandom\ on distinct properties violated per
budget, and to reach the first violation in fewer runs, because guard mutation targets the
conditions properties depend on whereas tool-sequence novelty rewards any behavioral
difference---most of which, in correctness testing, is benign variation. \todo{write to
the measured result; do not overclaim if the gap is small on some skills. Report per-skill
so mixed outcomes are visible.}

\subsection{RQ2: Component agreement and property strength}
\label{sec:rq2}

The model-driven components are calibrated against human labels: segmentation boundary
agreement (\SegF\ F1 on \ph{$\sim$100} hand-segmented traces), merge-decision agreement on
a labeled episode-pair set (\MergeAgree), and compiled-check agreement with hand-checked
verdicts (\CompileAgree). For property strength, we inject known violations---deleting a
test, weakening assertions to a trivial pass, declaring victory without running the suite,
force-pushing, pathological repetition---and report the detection rate. In a preliminary
run on one skill the suite detected \ph{all five} injected violation classes; we report
the full rate across the suite in \cref{tab:detect}. \todo{scale the injection study.}

\begin{table}[t]
\centering\small
\caption{Injected-violation detection. \todo{fill across suite.}}
\label{tab:detect}
\begin{tabular}{@{}l c@{}}
\toprule
\textbf{Injected violation} & \textbf{Detected}\\
\midrule
delete target test        & \ph{$\checkmark$}\\
weaken test to trivial pass & \ph{$\checkmark$}\\
declare victory, no test run & \ph{$\checkmark$}\\
force-push                & \ph{$\checkmark$}\\
pathological repetition   & \ph{$\checkmark$}\\
\midrule
Overall detection rate    & \DetectRate\\
\bottomrule
\end{tabular}
\end{table}

\subsection{RQ3: Improving an LLM-generated skill}
\label{sec:rq3}

Most circulating skills are LLM-generated. We test whether \tool's findings make one
better on a dedicated benchmark of \NumScenarios\ skill-generation \emph{scenarios}
(text/CSV/JSON processing, SQLite querying, validators, CLIs, config parsing, bug fixing,
algorithm implementation, log parsing). Each scenario ships \NumTestsPer\ \emph{held-out}
tests---each a prompt, an environment, and runnable pass-criteria---for \NumTestsTot\ tests
and \NumChecksTot\ executable checks in total, \emph{authored independently of any revision
loop} (by hand, disjoint from anything a tester generates and never shown to any tester or
reviser); this is the leakage control. A test passes iff all its checks pass; checks are
facet-driven (each test carries as many as it has real correctness facets). Every check was
verified in-container against a reference solution, so a test failure reflects the skill,
not a broken oracle. For each scenario we compare four versions of the \emph{one}
zero-shot LLM-generated base skill on the held-out tests: unrevised, and revised from the
Floor, Greybox, and \tool campaign reports respectively. The reviser prompt is
\emph{identical} across conditions; only the feedback payload differs. A difference in
held-out pass rate is therefore attributable to feedback quality, i.e.\ to the testing
method.

\begin{table}[t]
\centering\small
\caption{Held-out pass rate of an LLM-generated skill, zero-shot vs.\ revised from each
method's findings. \todo{fill.}}
\label{tab:rq3}
\begin{tabular}{@{}l c@{}}
\toprule
\textbf{Skill version} & \textbf{Held-out pass rate}\\
\midrule
Zero-shot (LLM-generated)     & \PassZeroShot\\
Revised from Floor findings   & \PassRandom\\
Revised from Greybox findings & \PassGreybox\\
Revised from \tool findings   & \PassSkillrace\\
\bottomrule
\end{tabular}
\end{table}

\subsection{RQ4: Ablations}
\label{sec:rq4}

We ablate (i) \tool's model, swapping \JudgeModel\ for a stronger and a weaker model and
reporting how yield moves; (ii) property-guided selection, replacing it with
uniform-random frontier choice, to isolate the value of directed falsification; and (iii)
the greybox granularity parameter, reporting yield at all three tool-event granularities.
\todo{fill from ablation runs.}

% =====================================================================
\section{Discussion and Threats to Validity}
\label{sec:threats}

\paragraph{What a \tool bug is, and is not.} A \tool report is evidence of a real defect:
a run that violated a property fixed in advance, with a replayable repro. It is not a
completeness claim---absence of violations does not certify a skill, and yield is bounded
by property strength (\cref{sec:props}). We treat this honestly by measuring detection of
planted violations rather than asserting coverage.

\paragraph{Model-driven components.} Segmentation, merging, guard extraction, and check
compilation are model judgments and can err. We contain this three ways: temperature-0
caching for within-campaign consistency; calibration against human labels (RQ2); and
architecture---a wrong merge surfaces as a downstream branch rather than corrupting a
verdict, and no model grades the run it produced. Build stability across selection order
is reported as an empirical number, not assumed.

\paragraph{Cost and parallelism.} The agent under test dominates cost; the validator
exists precisely to not waste it. \tool is more sequential than the baselines---each run
updates the tree that chooses the next input---so we report both per-run yield (where the
mechanism should win) and wall-clock (where the baselines parallelize more freely), rather
than a single number that hides the trade-off.

\paragraph{Generality.} Our suite is coding skills with mechanically checkable artifacts;
skills whose success is aesthetic or subjective are out of scope, and we mark the one
weakly-environment-contingent skill in our suite as such rather than averaging it into the
headline. The approach assumes the agent emits reasoning; agents that do not would degrade
\tool toward the reasoning-blind baselines.

% =====================================================================
\section{Related Work}
\label{sec:related}

\paragraph{Concolic and search-based testing.} Symbolic execution~\cite{king1976symbolic}
and its concrete-symbolic descendants---DART~\cite{dart}, CUTE~\cite{sen2005cute},
KLEE~\cite{klee}, and whitebox fuzzers such as SAGE~\cite{godefroid2008sage}---explore a
program's path space by reading branch conditions off the source and negating them to
reach new paths, using the path tree as both product and search signal. \tool imports this
loop but replaces its premise: an agent has no source, so we read candidate branch
conditions from the agent's reasoning and confirm or refute them by mutation. Recent work
brings concolic ideas to agents in a different direction: ConcoLLMic~\cite{concollmic} uses
an agent to \emph{perform} concolic execution on programs; we use concolic \emph{structure}
to test the agent's guidance. Coverage-guided greybox fuzzing~\cite{afl,ossfuzz} and
property-based testing~\cite{claessen2000quickcheck} supply our floor and our oracle
philosophy respectively: a property is an assertion fixed in advance, and a violation is a
real bug.

\paragraph{Testing LLM agents.} Agentic behavior is hard to test because it is
nondeterministic and mediated by tool choices~\cite{verigrey}. Most work targets
\emph{security}: AgentDojo~\cite{debenedetti2024agentdojo} benchmarks indirect prompt
injection~\cite{greshake2023not,owasp2025llm01,liu2025youraishell}, and
VeriGrey~\cite{verigrey} applies greybox fuzzing to discover injection prompts, using the
novelty of the tool-invocation sequence as feedback. We adopt VeriGrey's feedback mechanism
as our greybox rung, but retarget it from injection to correctness and hold the oracle
fixed, so the comparison speaks to feedback for correctness testing rather than to
VeriGrey's security performance. Defenses that constrain the agent rather than test
it---CaMeL~\cite{debenedetti2025camel}, instruction hierarchies~\cite{wallace2024instructionhierarchy},
and firewalls whose benchmarks may be saturated by simple filters~\cite{bhagwatkar2025firewalls}---%
are complementary: they harden deployment, whereas \tool tests the guidance an agent
follows. Our focus is orthogonal to injection: a skill can be perfectly safe and still be
\emph{wrong}, sending an agent confidently down a path its guidance does not cover.

\paragraph{Skills as a software artifact.} Skills are a recent and fast-spreading
convention~\cite{anthropic2025skills,anthropic2025skillseng}, distributed through shared
directories and, like rule files, a documented attack surface~\cite{pillar2025rules}; the
Model Context Protocol similarly externalizes agent capability~\cite{hou2025mcp}. Skills are
tested today the way early scripts were: by running them and seeing if they seem to work.
Benchmarks such as SWE-bench~\cite{jimenez2024swebench} evaluate whether an agent resolves a
task, but not whether a \emph{skill}'s guidance is robust across the situations an agent will
meet. \tool addresses that gap directly, and---because most skills are themselves
LLM-generated---closes the loop by using its findings to improve them.

% =====================================================================
\section{Conclusion}
\label{sec:conclusion}

Skills have become a unit of reuse for coding agents, but they are prose without a test
suite, trusted on the strength of the situations their authors happened to try. \tool
brings the discipline of concolic execution to this setting: it treats episodes as basic
blocks and the agent's stated reasons as branch conditions, mutates those conditions to
synthesize inputs that drive a skill down not-yet-taken branches, validates each input
before spending an agent run on it, and judges correctness by properties fixed in advance
rather than by a model's opinion of the run. \todo{close with the measured headline once
experiments land, and the RQ3 result on improving generated skills.}

\section*{Data Availability}
All code, the skill suite, generated inputs, and analysis scripts will be released.
\todo{artifact link.}

\bibliography{refs}
\bibliographystyle{plain}

\end{document}
